\documentclass{amsart}
% For dealing with references we use the comment environment
\usepackage{verbatim}
\newenvironment{reference}{\comment}{\endcomment}
%\newenvironment{reference}{}{}
\newenvironment{slogan}{\comment}{\endcomment}
\newenvironment{history}{\comment}{\endcomment}

% For commutative diagrams we use Xy-pic
\usepackage[all]{xy}

% We use 2cell for 2-commutative diagrams.
\xyoption{2cell}
\UseAllTwocells

% We use multicol for the list of chapters between chapters
\usepackage{multicol}

% This is generally recommended for better output
\usepackage{lmodern}
\usepackage[T1]{fontenc}

% For cross-file-references

% Package for hypertext links:
\usepackage{hyperref}

% For any local file, say "hello.tex" you want to link to please

% Theorem environments.
%
\theoremstyle{plain}
\newtheorem{theorem}[subsection]{Theorem}
\newtheorem{proposition}[subsection]{Proposition}
\newtheorem{lemma}[subsection]{Lemma}

\theoremstyle{definition}
\newtheorem{definition}[subsection]{Definition}
\newtheorem{example}[subsection]{Example}
\newtheorem{exercise}[subsection]{Exercise}
\newtheorem{situation}[subsection]{Situation}

\theoremstyle{remark}
\newtheorem{remark}[subsection]{Remark}
\newtheorem{remarks}[subsection]{Remarks}

\numberwithin{equation}{subsection}

% Macros
%
\def\lim{\mathop{\mathrm{lim}}\nolimits}
\def\colim{\mathop{\mathrm{colim}}\nolimits}
\def\Spec{\mathop{\mathrm{Spec}}}
\def\Hom{\mathop{\mathrm{Hom}}\nolimits}
\def\Ext{\mathop{\mathrm{Ext}}\nolimits}
\def\SheafHom{\mathop{\mathcal{H}\!\mathit{om}}\nolimits}
\def\SheafExt{\mathop{\mathcal{E}\!\mathit{xt}}\nolimits}
\def\Sch{\mathit{Sch}}
\def\Mor{\mathop{\mathrm{Mor}}\nolimits}
\def\Ob{\mathop{\mathrm{Ob}}\nolimits}
\def\Sh{\mathop{\mathit{Sh}}\nolimits}
\def\NL{\mathop{N\!L}\nolimits}
\def\CH{\mathop{\mathrm{CH}}\nolimits}
\def\proetale{{pro\text{-}\acute{e}tale}}
\def\etale{{\acute{e}tale}}
\def\QCoh{\mathit{QCoh}}
\def\Ker{\mathop{\mathrm{Ker}}}
\def\Im{\mathop{\mathrm{Im}}}
\def\Coker{\mathop{\mathrm{Coker}}}
\def\Coim{\mathop{\mathrm{Coim}}}

% Boxtimes
%
\DeclareMathSymbol{\boxtimes}{\mathbin}{AMSa}{"02}

%
% Macros for moduli stacks/spaces
%
\def\QCohstack{\mathcal{QC}\!\mathit{oh}}
\def\Cohstack{\mathcal{C}\!\mathit{oh}}
\def\Spacesstack{\mathcal{S}\!\mathit{paces}}
\def\Quotfunctor{\mathrm{Quot}}
\def\Hilbfunctor{\mathrm{Hilb}}
\def\Curvesstack{\mathcal{C}\!\mathit{urves}}
\def\Polarizedstack{\mathcal{P}\!\mathit{olarized}}
\def\Complexesstack{\mathcal{C}\!\mathit{omplexes}}
% \Pic is the operator that assigns to X its picard group, usage \Pic(X)
% \Picardstack_{X/B} denotes the Picard stack of X over B
% \Picardfunctor_{X/B} denotes the Picard functor of X over B
\def\Pic{\mathop{\mathrm{Pic}}\nolimits}
\def\Picardstack{\mathcal{P}\!\mathit{ic}}
\def\Picardfunctor{\mathrm{Pic}}
\def\Deformationcategory{\mathcal{D}\!\mathit{ef}}

\setlength{\emergencystretch}{3em}
\begin{document}
\title[Stacks Project, Tag 094B]{663 --- Compact derived objects are retracts of finite generator complexes}
\maketitle
\noindent Copyright (C) 2005--2025 Johan de Jong. Stacks Project authors. Source: \href{https://stacks.math.columbia.edu/tag/094B}{Tag 094B}.

\noindent Editorial difficulty note (not author text). Difficulty H2: The proof first obtains a bounded retract via truncation homotopy colimits, then inductively shrinks each infinite sum using compactness and homological lifting through a distinguished triangle. This finite reconstruction argument is a substantial compact-generation theorem..

\noindent Editorial provenance note (not author text). History: excerpt from revision \texttt{a0444\allowbreak 6e57e\allowbreak c1fbc\allowbreak 252a8\allowbreak 71afc\allowbreak ec775\allowbreak 2fb28\allowbreak 07b14}, sites-cohomology.tex, lines 13933--14032. Reference presentation was adapted; original-author context and core support blocks were retained or added. The mathematical wording is unchanged. Licensed under GNU FDL 1.2 or later; no invariant sections or cover texts. See ../licenses/COPYING. Context blocks reproduce author definitions and supporting results; they are not additional counted problems.

\begin{definition}
\label{derived-definition-compact-object}
Let $\mathcal{D}$ be an additive category with arbitrary direct
sums. A {\it compact object} of $\mathcal{D}$ is an object $K$
such that the map
$$
\bigoplus\nolimits_{i \in I} \Hom_{\mathcal{D}}(K, E_i)
\longrightarrow
\Hom_{\mathcal{D}}(K, \bigoplus\nolimits_{i \in I} E_i)
$$
is bijective for any set $I$ and objects
$E_i \in \Ob(\mathcal{D})$ parametrized by $i \in I$.
\end{definition}

\begin{lemma}
\label{lemma-compact-in-terms-of-generators}
Let $\mathcal{A}$ be a Grothendieck abelian category. Let
$S \subset \Ob(\mathcal{A})$ be a set of objects such that
\begin{enumerate}
\item any object of $\mathcal{A}$ is a quotient of a direct sum
of elements of $S$, and
\item for any $E \in S$ the functor $\Hom_\mathcal{A}(E, -)$
commutes with direct sums.
\end{enumerate}
Then every compact object of $D(\mathcal{A})$ is a direct summand
in $D(\mathcal{A})$ of a finite complex of finite direct sums of
elements of $S$.
\end{lemma}

\begin{proof}
Assume $K \in D(\mathcal{A})$ is a compact object. Represent $K$ by a complex
$K^\bullet$ and consider the map
$$
K^\bullet
\longrightarrow
\bigoplus\nolimits_{n \geq 0} \tau_{\geq n} K^\bullet
$$
where we have used the canonical truncations, see
Homology, Section \href{https://stacks.math.columbia.edu/tag/0118}{section truncations (Tag 0118)}.
This makes sense as in each degree the direct sum on the right is finite.
By assumption this map factors through a finite direct sum.
We conclude that $K \to \tau_{\geq n} K$ is zero for at least one $n$,
i.e., $K$ is in $D^{-}(R)$.

\medskip\noindent
We may represent $K$ by a bounded above complex $K^\bullet$ each of whose
terms is a direct sum of objects from $S$, see
Derived Categories, Lemma \href{https://stacks.math.columbia.edu/tag/05T7}{lemma subcategory left resolution (Tag 05T7)}.
Note that we have
$$
K^\bullet = \bigcup\nolimits_{n \leq 0} \sigma_{\geq n}K^\bullet
$$
where we have used the stupid truncations, see
Homology, Section \href{https://stacks.math.columbia.edu/tag/0118}{section truncations (Tag 0118)}.
Hence by Derived Categories, Lemmas \href{https://stacks.math.columbia.edu/tag/0949}{lemma colim hocolim (Tag 0949)} and
\href{https://stacks.math.columbia.edu/tag/094A}{lemma commutes with countable sums (Tag 094A)}
we see that $1 : K^\bullet \to K^\bullet$ factors through
$\sigma_{\geq n}K^\bullet \to K^\bullet$ in $D(R)$.
Thus we see that $1 : K^\bullet \to K^\bullet$ factors as
$$
K^\bullet \xrightarrow{\varphi} L^\bullet \xrightarrow{\psi} K^\bullet
$$
in $D(\mathcal{A})$ for some complex $L^\bullet$ which is bounded and
whose terms are direct sums of elements of $S$. Say $L^i$ is zero for
$i \not \in [a, b]$. Let $c$ be the largest integer $\leq b + 1$ such
that $L^i$ a finite direct sum of elements of $S$ for $i < c$.
Claim: if $c < b + 1$, then we can modify $L^\bullet$ to increase $c$.
By induction this claim will show we have a factorization
of $1_K$ as
$$
K \xrightarrow{\varphi} L \xrightarrow{\psi} K
$$
in $D(\mathcal{A})$ where $L$ can be represented by a finite
complex of finite direct sums of elements of $S$. Note that
$e = \varphi \circ \psi \in \text{End}_{D(\mathcal{A})}(L)$
is an idempotent. By Derived Categories,
Lemma \href{https://stacks.math.columbia.edu/tag/05QW}{lemma projectors have images triangulated (Tag 05QW)}
we see that $L = \Ker(e) \oplus \Ker(1 - e)$.
The map $\varphi : K \to L$ induces an isomorphism with
$\Ker(1 - e)$ in $D(R)$ and we conclude.

\medskip\noindent
Proof of the claim. Write $L^c = \bigoplus_{\lambda \in \Lambda} E_\lambda$.
Since $L^{c - 1}$ is a finite direct sum of elements of $S$
we can by assumption (2) find a finite subset
$\Lambda' \subset \Lambda$ such that $L^{c - 1} \to L^c$ factors
through $\bigoplus_{\lambda \in \Lambda'} E_\lambda \subset L^c$.
Consider the map of complexes
$$
\pi :
L^\bullet
\longrightarrow
(\bigoplus\nolimits_{\lambda \in \Lambda \setminus \Lambda'} E_\lambda)[-i]
$$
given by the projection onto the factors corresponding to
$\Lambda \setminus \Lambda'$ in degree $i$.
By our assumption on $K$ we see that, after possibly replacing $\Lambda'$ by
a larger finite subset, we may assume that $\pi \circ \varphi = 0$
in $D(\mathcal{A})$.
Let $(L')^\bullet \subset L^\bullet$ be the kernel of $\pi$.
Since $\pi$ is surjective we get a short exact sequence of complexes,
which gives a distinguished triangle in $D(\mathcal{A})$ (see
Derived Categories, Lemma \href{https://stacks.math.columbia.edu/tag/0152}{lemma derived canonical delta functor (Tag 0152)}).
Since $\Hom_{D(\mathcal{A})}(K, -)$ is homological (see
Derived Categories, Lemma \href{https://stacks.math.columbia.edu/tag/0149}{lemma representable homological (Tag 0149)})
and $\pi \circ \varphi = 0$, we can find a morphism
$\varphi' : K^\bullet \to (L')^\bullet$ in $D(\mathcal{A})$ whose
composition with $(L')^\bullet \to L^\bullet$ gives $\varphi$.
Setting $\psi'$ equal to the composition of $\psi$ with
$(L')^\bullet \to L^\bullet$ we obtain a new factorization.
Since $(L')^\bullet$ agrees with $L^\bullet$ except in degree $c$
and since $(L')^c = \bigoplus_{\lambda \in \Lambda'} E_\lambda$ the
claim is proved.
\end{proof}
\end{document}
